Une patineuse tourne sur elle-même à \qty{\v{f1}}{tr/s}, les bras tendus ; son moment d'inertie est alors de \qty{\v{I1}}{kg.m^2}. Elle ramène les bras contre son corps et son moment d'inertie devient \qty{\v{I2}}{kg.m^2}. On néglige le frottement de la glace.

\begin{enumerate}[label=\alph*)]
    \item Déterminez sa vitesse angulaire finale, en radians par seconde.
    \item Déterminez son énergie cinétique de rotation initiale.
    \item Déterminez la variation de son énergie cinétique de rotation. D'où vient cette énergie ?
\end{enumerate}
